\section{Counting calls and bit operations}
\label{sec:complexity}

Correctness now gives the complete winning set. To finish the theorem,
we must bound the work in terms of binary input length rather than the
numerical width of the box. There are two relevant depths: the width
can be halved polynomially many times in $L$, but the mass budget can
decrease only logarithmically many times in $n$. Because only one
child preserves that budget, these depths give a quasipolynomial call
count. We then bound the exact arithmetic within each call.

\subsection{The recursion tree}

For the initial data of Equation~\eqref{eq:main-call}, write
\begin{equation}
 d_0=\log_2 D_0,\qquad B_0=B(1/n^2),\qquad
 H=d_0+B_0,\qquad s=2048n+3.
 \label{eq:initial-resources}
\end{equation}
The minimum initial mass product is $1/n^2$. Every path in the call
tree has at most $H$ edges by the rank decrease of
Lemma~\ref{lem:termination}. At most $B_0$ of its edges can
decrease $B$, even if an individual edge decreases it by more than
one. At any node there is at most one child preserving $B$,
namely the provisional child. There are at most $2(1024n+1)$ loop
children and one continuation, so at most $s$ children decrease $B$.

\begin{lemma}[Number of calls]
\label{lem:node-count}
The complete main computation makes at most
\begin{equation}
 \sum_{\ell=0}^{H}\sum_{j=0}^{\min(\ell,B_0)}
       \binom{\ell}{j}s^j
 \ \le\ (H+1)(B_0+1)(1+Hs)^{B_0}
 \label{eq:node-bound}
\end{equation}
calls, including the initial call.
\end{lemma}

\begin{proof}
Consider a node at depth $\ell$ whose path uses $j$ budget-decreasing
edges. Their positions can be chosen in $\binom{\ell}{j}$ ways.
Order each node's budget-decreasing children by execution order.
At each chosen position there are at most $s$ choices of that ordinal; at each
other position there is at most one. This specifies at most
$\binom{\ell}{j}s^j$ possible nodes. Sum over all depths
$0\le\ell\le H$ and $0\le j\le\min(\ell,B_0)$.
For the second inequality, each summand is at most
$(Hs)^j\le(1+Hs)^{B_0}$ and there are at most
$(H+1)(B_0+1)$ summands.
\end{proof}

This count uses the same recursion principle as the reduced-precision
analysis of Parys~\cite[Section~5]{Parys2019}, further developed in
\cite{LehtinenParysScheweWojtczak2022}: one child retains the bounded
parameter while all others improve it. Here the two parameters are
the logarithm of an integer width and the budget derived from rational
mass products.

The explicit encoding bounds both $n$ and $m=|E|$ polynomially in
$L+1$, and each original weight has $O(L+1)$ bits. The perturbation
and Equation~\eqref{eq:initial-width} imply
\[
 \log_2(W+1)=O(L+\log(n+1)),\qquad
 d_0=O(L+\log(n+1)).
\]
Since $\rho=32/25$ is a fixed constant greater than one,
\[
 B_0=1+\lfloor\log_{\rho}(n^2)\rfloor
                          =O(1+\log(n+1)).
\]
Thus $H$ and $s$ are polynomial in $L+1$, and
$B_0=O(\log(L+2))$. Taking logarithms in the call bound gives
\begin{equation}
 \begin{split}
 \log_2\bigl(\text{number of calls}\bigr)
 &\le \log_2(H+1)+\log_2(B_0+1)
                         +B_0\log_2(1+Hs)\\
 &=O((\log_2(L+2))^2).
 \end{split}
 \label{eq:quasipoly-calls}
\end{equation}
This includes $n=1$, for which $B_0=1$.

\subsection{Exact representations}

The nested-box property keeps every endpoint and pivot of the main
computation inside $[0,D_0]^V$. Each such coordinate and each width
therefore uses $O(d_0+1)$ bits. In a basic iteration, an intermediate
edge expression is a coordinate plus a transformed weight. It may be
negative, but its signed binary length is
$O(d_0+\log(W+1)+1)$, polynomial in $L+1$.

Represent a mass by a positive numerator and a positive denominator;
reducing fractions is unnecessary. Starting from $1/n$, a mass is
multiplied along any path only by numbers in
\[
                           \{1,4/3,4/5,8/5\}.
\]
Only budget-decreasing edges change the masses, and a path contains
at most $B_0$ such edges. Thus each numerator and denominator has
$O(B_0+\log(n+1))$ bits. Children receive their own copies of the
input masses. In particular, completing one sibling does not multiply
the parent masses before another sibling starts. Testing $a_i b_i>1$
uses integer cross-multiplication, and testing $b_i>1$ or applying a
constant multiplier also uses numbers of polynomial length. All mass
arithmetic is exact.

The proof quantities $B$ and $\mu$, and the peaks, gaps, supports,
and witness vectors, need not be stored or computed. The algorithm
uses neither a logarithm evaluation for its rank nor an oracle for
the existence of a witness. The procedure evaluates $F$ only inside
$\Basic$, by explicit max/min scans of the edge list.

\subsection{Local cost and uniform implementation}

Each call has polynomial nonrecursive work. The mass and width tests
are finite scans with the arithmetic just bounded. If the width base
case applies, $D<4096$, and the imported iteration bound gives
\[
                            1+nD\le1+4095n
\]
evaluations, including the final stability test. Every evaluation
scans the $m$ explicit edges, adds their transformed weights, takes
the relevant extrema, and clips at $n$ endpoints. It acts on integers
of polynomial length. The midpoint comparison is implemented as
$2z_i>2A_i+D$, which also handles $D=1$.

In a nonbase call, the constants defining the pivot are obtained by
exact divisions by powers of two. Each pass has $O(n+1)$ iterations
apart from its children's work. A scan of the $n$ coordinates performs
one simultaneous update and tests whether it changed anything.
Preparing child inputs and receiving their labels, forming $I_-$
and $I_+$, reweighting, and applying the final overrides likewise
use polynomially many operations on polynomial-length integers.
Schoolbook arithmetic suffices throughout.

A depth-first implementation runs children sequentially and retains
one record per active call. The record holds the parent masses,
endpoints, pivot, labels, and current loop position. There are at most
$H+1$ active records, each of polynomial size. Stack management is
therefore included in the same polynomial local bound. Preprocessing
the transformed weights, finding the least power of two in
Equation~\eqref{eq:initial-width}, and creating masses $1/n$ take
polynomial time. Reading the input and producing the $n$ final labels
also do so.

Multiply this polynomial per-call cost by the call count bounded in
Equation~\eqref{eq:quasipoly-calls}, and include preprocessing and
output. For an absolute constant $C>0$, the complete bound is
\begin{equation}
                            2^{C(\log_2(L+2))^2}.
 \label{eq:final-bit-bound}
\end{equation}
The constants in the procedure are fixed independently of the game,
and every arithmetic and control operation has been specified. This
gives a uniform deterministic bit algorithm.

\begin{proof}[Proof of Theorem~\ref{thm:main}]
Compute Equation~\eqref{eq:main-call}, run $\Label$, and output
the positively labelled vertices. Theorem~\ref{thm:labelling}
guarantees the labels required by every witness. In the initial box,
the box solution is a witness of both types, and the argument after
Equation~\eqref{eq:main-call} identifies all its required labels
with the bands of Proposition~\ref{prop:game-reduction}. The
positive band is exactly the zero-threshold winning set, including
against every history-dependent opponent. Equation~\eqref{eq:final-bit-bound}
bounds the entire computation. Signed binary weights, self-loops,
and separately listed parallel edges are all covered by the
construction and accounting above.
\end{proof}
